\begin{thebibliography}{99}

\bibitem{esbensen2010proper}
Esbensen KH, Geladi P. Principles of proper validation: use and abuse of re-sampling for validation. \textit{J. Chemom.} 2010;24(3--4):168--187. doi:10.1002/cem.1310.

\bibitem{tropsha2003earnest}
Tropsha A, Gramatica P, Gombar VK. The importance of being earnest: validation is the absolute essential for successful application and interpretation of QSPR models. \textit{QSAR Comb. Sci.} 2003;22(1):69--77. doi:10.1002/qsar.200390007.

\bibitem{gramatica2007principles}
Gramatica P. Principles of QSAR models validation: internal and external. \textit{QSAR Comb. Sci.} 2007;26(5):694--701. doi:10.1002/qsar.200610151.

\bibitem{lasfar2024robustness}
Lasfar R, T\'{o}th G. The difference of model robustness assessment using cross-validation and bootstrap methods. \textit{J. Chemom.} 2024;38(6):e3530. doi:10.1002/cem.3530.

\bibitem{kiraly2025leakage}
Kir\'{a}ly P, T\'{o}th G. Being aware of data leakage and cross-validation scaling in chemometric model validation. \textit{J. Chemom.} 2025;39(4):e70026. doi:10.1002/cem.70026.

\bibitem{camacho2026validation}
Camacho J. A set of rules for model validation. \textit{J. Chemom.} 2026;40(2):e70110. doi:10.1002/cem.70110.

\bibitem{wu2018moleculenet}
Wu Z, Ramsundar B, Feinberg EN, et al. MoleculeNet: a benchmark for molecular machine learning. \textit{Chem. Sci.} 2018;9(2):513--530. doi:10.1039/C7SC02664A.

\bibitem{bemis1996properties}
Bemis GW, Murcko MA. The properties of known drugs. 1. Molecular frameworks. \textit{J. Med. Chem.} 1996;39(15):2887--2893. doi:10.1021/jm9602928.

\bibitem{yang2019learned}
Yang K, Swanson K, Jin W, et al. Analyzing learned molecular representations for property prediction. \textit{J. Chem. Inf. Model.} 2019;59(8):3370--3388. doi:10.1021/acs.jcim.9b00237.

\bibitem{deng2023systematic}
Deng J, Yang Z, Wang H, Ojima I, Samaras D, Wang F. A systematic study of key elements underlying molecular property prediction. \textit{Nat. Commun.} 2023;14:6395. doi:10.1038/s41467-023-41948-6.

\bibitem{joeres2025datasail}
Joeres R, Blumenthal DB, Kalinina OV. Data splitting to avoid information leakage with DataSAIL. \textit{Nat. Commun.} 2025;16:3337. doi:10.1038/s41467-025-58606-8.

\bibitem{landrum2023simpd}
Landrum GA, Beckers M, Lanini J, Schneider N, Stiefl N, Riniker S. SIMPD: an algorithm for generating simulated time splits for validating machine learning approaches. \textit{J. Cheminform.} 2023;15:119. doi:10.1186/s13321-023-00787-9.

\bibitem{li2026partition}
Li S, Sun P. Toward more trustworthy QSAR: a systematic discussion on data set partitioning. \textit{J. Chem. Inf. Model.} 2026;66(4):2199--2210. doi:10.1021/acs.jcim.5c02465.

\bibitem{nael2026dataset}
Nael MA, Alakonda LM, Elokely KM. Defining the data set defines the QSAR claim. \textit{J. Chem. Inf. Model.} 2026;66(6):2951--2954. doi:10.1021/acs.jcim.6c00514.

\bibitem{netzeva2005applicability}
Netzeva TI, Worth A, Aldenberg T, et al. Current status of methods for defining the applicability domain of (quantitative) structure-activity relationships: the report and recommendations of ECVAM Workshop 52. \textit{Altern. Lab. Anim.} 2005;33(2):155--173. doi:10.1177/026119290503300209.

\bibitem{sahigara2012applicability}
Sahigara F, Mansouri K, Ballabio D, Mauri A, Consonni V, Todeschini R. Comparison of different approaches to define the applicability domain of QSAR models. \textit{Molecules.} 2012;17(5):4791--4810. doi:10.3390/molecules17054791.

\bibitem{ezenarro2025validation}
Ezenarro J, Schorn-Garc\'{i}a D. How are chemometric models validated? A systematic review of linear regression models for NIRS data in food analysis. \textit{J. Chemom.} 2025;39(6):e70036. doi:10.1002/cem.70036.

\bibitem{vantilborg2022cliffs}
van Tilborg D, Alenicheva A, Grisoni F. Exposing the limitations of molecular machine learning with activity cliffs. \textit{J. Chem. Inf. Model.} 2022;62(23):5938--5951. doi:10.1021/acs.jcim.2c01073.

\bibitem{jimenez2022attribution}
Jim\'{e}nez-Luna J, Skalic M, Weskamp N. Benchmarking molecular feature attribution methods with activity cliffs. \textit{J. Chem. Inf. Model.} 2022;62(2):274--283. doi:10.1021/acs.jcim.1c01163.

\bibitem{rdkit}
Landrum G, Tosco P, Kelley B, et al. RDKit: open-source cheminformatics, Release 2026.03.4. Zenodo. 2026. doi:10.5281/zenodo.21291217.

\bibitem{rogers2010extended}
Rogers D, Hahn M. Extended-connectivity fingerprints. \textit{J. Chem. Inf. Model.} 2010;50(5):742--754. doi:10.1021/ci100050t.

\bibitem{pedregosa2011scikit}
Pedregosa F, Varoquaux G, Gramfort A, et al. Scikit-learn: machine learning in Python. \textit{J. Mach. Learn. Res.} 2011;12:2825--2830.

\bibitem{breiman2001random}
Breiman L. Random forests. \textit{Mach. Learn.} 2001;45:5--32. doi:10.1023/A:1010933404324.

\bibitem{chen2016xgboost}
Chen T, Guestrin C. XGBoost: a scalable tree boosting system. In: \textit{Proceedings of the 22nd ACM SIGKDD International Conference on Knowledge Discovery and Data Mining}; 2016:785--794. doi:10.1145/2939672.2939785.

\end{thebibliography}
